\documentclass[spanish,oneside]{book} % Este documento es en español
%\documentclass[spanish,journal]{IEEEtran} % En caso de que me exijan formato IEEE
\def\spanishoptions{mexico}
\usepackage[spanish]{babel}      % Para procesar adecuadamente el español
\usepackage[utf8]{inputenc}      % Codificación de entrada UTF-8
\usepackage[T1]{fontenc}         % Codificación de salida T1
\usepackage[official]{eurosym}   % Símbolo de euro
\usepackage{siunitx}             % Paquete para unidades métricas
\usepackage{textcomp}            % Tipografías compañeras
\usepackage{lmodern}             % Tipografía Latin Modern
\usepackage{moreverb}            % Para escribir código
\usepackage{tfrupee}
\usepackage{fontawesome5}
\usepackage{parskip}
\usepackage{hyperref}
\usepackage{tasks}
\usepackage{exsheets}
\usepackage{enumitem}
\usepackage{lastpage}
\usepackage{booktabs}
\usepackage{graphicx}
\usepackage{apalike}
\usepackage{relsize}
%\usepackage{inlinebib}
%\usepackage{mathptmx}
%\usepackage{palatino}
%\usepackage{helvet}

% Formato global:
\setlength{\parindent}{0cm}      % Párrafos sin sangría inicial
\hypersetup{pdfborder={0 0 0}}   % Ligas sin borde

\begin {document}
\title {Guía de estudio de la clase de Aprendizaje Machín}
\author {Acoyani Garrido Sandoval}
\date {Septiembre de 2026}
\maketitle

\tableofcontents


\chapter {Conceptos previos}

\section {Métricas MSE y R2}




% \bibliographystyle{apalike}
% \bibliographystyle{inlinebib}
% \bibliography{estudiogeneral}
\end{document}
